% =========================
% Section 3: Data Processing and EDA (model-driven)
% =========================


\subsection{Data Cleaning Principles and Validation}
\label{subsec:cleaning_principles_validation}

\subsubsection{Cleaning Principles}
\label{subsubsec:cleaning_principles}

Our preprocessing is \emph{model-driven}: we only perform transformations required to
(i) construct a reproducible season--week--contestant panel,
(ii) define eligibility/denominator rules for KPI computation, and
(iii) enable counterfactual replay under alternative voting schemes.
We preserve the dataset's official encodings and remove their influence through an explicit
\emph{active-set} definition rather than aggressive imputation \cite{comap2026_problemC_pdf}.

\paragraph{Principle A --- Respect official encodings (0 scores, N/A, irregular season length).}
The official problem statement clarifies that (i) weeks with no elimination or multiple eliminations exist,
(ii) \texttt{N/A} can mean either ``no 4th judge'' or ``weeks beyond a season's actual run'',
and (iii) a score of 0 is recorded for celebrities after elimination (post-exit padding).
We therefore treat these as \emph{structural encodings} rather than noise \cite{comap2026_problemC_pdf}.

\paragraph{Principle B --- Eligibility controls denominators (active-set construction).}
A contestant is considered active in $(s,t)$ if at least one judge score exists for that week and the week
does not exceed the contestant's exit week (eliminated week parsed from \texttt{results}, or an inferred
withdraw week). Post-elimination zeros are retained for auditability but excluded from denominators via
\texttt{active=\textbf{false}}.

\paragraph{Principle C --- Within-week aggregation only; no judge identity tracking.}
Because judge indices are not persistent identities across weeks/seasons, we only use within-week aggregates
(weekly total, rank, and percent) and avoid modeling judge-specific time series at the preprocessing stage
\cite{comap2026_problemC_pdf}.

\paragraph{Principle D --- Non-single-elimination weeks are handled explicitly.}
Weeks with zero eliminations are excluded from the single-elimination consistency denominator, while
multi-elimination weeks are evaluated via bottom-$k$ set matching (set-match).


\subsubsection{Data Cleaning Operations}
\label{subsubsec:cleaning_operations}

We implement a deterministic pipeline that reshapes the original wide-format judge-score columns into a long
panel and then attaches minimal labels required for downstream inference and replay:

\begin{enumerate}
  \item \textbf{Wide-to-long reshaping.}
  We reshape \texttt{weekX\_judgeY\_score} into one row per $(season, week, celebrity)$, and compute
  \texttt{all\_judges\_nan} and
  \texttt{total\_judge\_score = sum(judge\_scores, skipna=\{True\})}
  to respect varying judge counts.

  \item \textbf{Exit mapping and labels.}
  We parse elimination week from \texttt{results} when available; otherwise, for withdrew/finished cases we infer
  \texttt{exit\_week} as the last week with a positive total score, falling back to the last non-missing scored week.
  We store \texttt{exit\_type} $\in \{\text{withdrew}, \text{eliminated}, \text{finished}\}$ and define
  \texttt{true\_elim\_flag} only for elimination exits by default.

  \item \textbf{Active-set definition and anomaly flag.}
  We define \texttt{active = (not all\_judges\_nan) \& (week $\le$ exit\_week)}.
  Additionally, if \texttt{total\_judge\_score == 0} occurs outside the inferred exit week, we flag
  \texttt{data\_anomaly\_zero\_score=1} and force \texttt{active=false} to prevent padded zeros from contaminating denominators.

  \item \textbf{Within-week relative measures.}
  For each $(season, week)$, among active contestants only, we compute \texttt{judge\_rank} (1 = best) and
  \texttt{judge\_percent = score / sum(scores)} when the denominator is positive.
\end{enumerate}


\subsubsection{Data Cleaning Outcomes and Validation}
\label{subsubsec:cleaning_outcomes_validation}

\paragraph{Outcomes (artifacts and interfaces).}
The output of preprocessing is a \emph{canonical weekly panel}
(\texttt{output/table/intermediate\_weekly\_panel.csv}) that contains eligibility/exit labels
(\texttt{active}, \texttt{exit\_type}, \texttt{exit\_week}, \texttt{true\_elim\_flag}) and judge-based within-week features
(\texttt{total\_judge\_score}, \texttt{judge\_rank}, \texttt{judge\_percent}).
This panel is the single input interface used for all downstream KPI recomputation, vote-share inference, and replay simulation.

In addition, we generate audit/debug artifacts: (i) a week-level prediction table
(\texttt{output/table/intermediate\_baseline\_preds.csv}) and (ii) a report-ready \LaTeX{} table
(\texttt{output/table/tab\_baseline\_consistency.tex}) plus an optional figure
(\texttt{output/figure/fig\_fliprate\_by\_season.pdf}).

\paragraph{Validation (what we check, and why it is sufficient for P0).}
We validate preprocessing through three audit layers:
\begin{itemize}
  \item \textbf{(V1) Encoding-consistency checks.}
  We follow the official semantics for \texttt{N/A} and post-elimination zeros, preserving these encodings and excluding them
  from denominators via \texttt{active} \cite{comap2026_problemC_pdf}.

  \item \textbf{(V2) Denominator/eligibility invariants.}
  \texttt{judge\_percent} is computed only among active contestants within each $(season, week)$; when total score is positive,
  it sums to one within the active set. The explicit \texttt{data\_anomaly\_zero\_score} flag prevents accidental inclusion of padded zeros.

  \item \textbf{(V3) End-to-end reproducibility via BL-0 sanity check.}
  On top of the canonical panel, we implement the BL-0 baseline (rank: judge-rank only; percent: judge-percent plus uniform
  $1/n_{\text{active}}$) and reproduce a consistent set of KPI1/KPI3 summaries across era slices.
  This serves as a pipeline sanity check rather than a final model claim.
\end{itemize}

\noindent\textbf{BL-0 sanity-check summary (reportable, not a final claim).}
Under BL-0, the overall elimination hit-rate is 0.364 (rank) and 0.375 (percent) across 264 eligible weeks,
and the rank-vs-percent flip rate is 0.038.


\subsection{Panel Construction}
\label{subsec:panel_construction}

Panel construction converts the wide-format score table into a long-format dataset where each record corresponds to a unique
$(season, week, celebrity)$ unit. We reshape \texttt{weekX\_judgeY\_score} into such a long panel and compute weekly aggregates
(\texttt{total\_judge\_score}) and within-week relative measures (\texttt{judge\_rank}, \texttt{judge\_percent}) among the active set.
We then attach deterministic exit/eligibility labels (\texttt{exit\_type}, \texttt{exit\_week}, \texttt{true\_elim\_flag}, \texttt{active})
derived from \texttt{results} and score availability. This panel is the canonical input shared by all subsequent modules, including
vote-share inference and counterfactual replay under alternative voting schemes.


\subsection{Empirical Findings from EDA (only what feeds models)}
\label{subsec:eda_model_driven}

\subsubsection{Data Visualization}
\label{subsubsec:eda_viz}

We summarize baseline consistency and rule divergence using:
\begin{itemize}
  \item \textbf{(L1 table)} \texttt{output/table/tab\_baseline\_consistency.tex}:
  BL-0 elimination hit-rate under rank and percent schemes (KPI1) and rule divergence (KPI3/FlipRate), with slices:
  Overall / Rank-era / Percent-era.
  \item \textbf{(L2 figure, optional)} \texttt{output/figure/fig\_fliprate\_by\_season.pdf}:
  season-level visualization of where the two schemes disagree.
\end{itemize}

% If you want to include the table/figure directly:
% \begin{table}[t]
%   \centering
%   \begin{tabular}{lrrrr}
\hline
Slice & $N_{\mathrm{weeks}}$ & Hit-Rate (Rank) & Hit-Rate (Percent) & FlipRate \\
\hline
Overall & 264 & 0.364 & 0.375 & 0.038 \\
Rank-era (S1-2,28-34) & 66 & 0.348 & 0.364 & 0.015 \\
Percent-era (S3-27) & 198 & 0.369 & 0.379 & 0.045 \\
\hline
\end{tabular}
%   \caption{BL-0 baseline consistency and rule divergence across era slices.}
%   \label{tab:baseline_consistency}
% \end{table}
%
% \begin{figure}[t]
%   \centering
%   \includegraphics[width=\linewidth]{output/figure/fig_fliprate_by_season.pdf}
%   \caption{Season-level flip rate between rank- and percent-based BL-0 predictions.}
%   \label{fig:fliprate_by_season}
% \end{figure}


\subsubsection{Key Empirical Findings}
\label{subsubsec:key_findings}

\paragraph{Finding 1 --- A reproducible BL-0 baseline exists and is non-trivial.}
Using only judge-based information and a uniform fan proxy, BL-0 yields measurable elimination-consistency hit-rates under both
rank and percent aggregation, and a non-zero flip rate between the two schemes. These quantities are used strictly as a sanity check
for the pipeline and as a reference point for later vote-inference models.

\paragraph{Finding 2 --- Rule comparisons must be stress-tested around the season-28 cutoff.}
We adopt the era split for rule-comparison slicing (Rank-era vs. Percent-era), but treat the season-28 switch as an assumption and
stress-test adjacent cutoffs.

\paragraph{Finding 3 --- Non-standard weeks materially affect denominators.}
Because the dataset includes weeks with no elimination or multiple eliminations, KPI denominators must explicitly exclude no-elimination
weeks and handle multi-elimination weeks via set matching; otherwise consistency metrics can be inflated.


\subsection{Feature Engineering}
\label{subsec:feature_engineering}

We keep feature engineering minimal and aligned with downstream modeling needs:

\begin{itemize}
  \item \textbf{Performance signals (week-level, within active set).}
  \texttt{total\_judge\_score}, \texttt{judge\_rank}, and \texttt{judge\_percent} are computed per $(season, week)$ among active contestants
  and serve as the observed performance component when replaying elimination rules or inferring latent vote shares.

  \item \textbf{Explanatory covariates (contestant-level, for heterogeneity analysis).}
  When available, we use \texttt{ballroom\_partner} (pro dancer identity), \texttt{celebrity\_age\_during\_season},
  \texttt{celebrity\_industry}, and region-related fields for slicing and mixed-effects heterogeneity analyses.
\end{itemize}


\subsection{Synthesis of Findings}
\label{subsec:synthesis_findings}

In summary, our preprocessing produces a canonical season--week--contestant panel that
(i) respects the official encodings of \texttt{N/A} and post-elimination zeros,
(ii) enforces denominator correctness through an explicit active-set definition, and
(iii) enables reproducible baseline replay consistent with our KPI definitions.
Key data caveats---non-standard elimination weeks, varying judge availability, and the uncertain season-28 era cutoff---are surfaced
explicitly for robustness stress tests rather than hidden by aggressive cleaning \cite{comap2026_problemC_pdf}.
